\chapter{Objectives}

Based on the problem statement and the scope achievable within the mini-project timeline, the following three objectives have been identified as the core, achievable deliverables of this project:

\begin{enumerate}
    \item \textbf{Build a unified knowledge graph} that links transaction, KYC, and behavioural data for a chosen set of accounts, so that relationships between accounts (shared devices, shared phone numbers, shared recipients) that are invisible in a purely tabular view become directly queryable.
    \item \textbf{Implement an AI-based transaction risk-scoring model} (gradient-boosted-tree and/or Graph-Neural-Network driven) that flags suspicious transactions in real time, using both the transaction's own features and context derived from the knowledge graph, together with an explanation for each score produced.
    \item \textbf{Automate generation of a compliant Suspicious Activity Report (SAR) draft} from flagged cases using a Retrieval-Augmented Generation (RAG) pipeline grounded in real regulatory text, reducing the manual drafting effort currently required of a compliance analyst.
\end{enumerate}

These three objectives were deliberately scoped to be independently achievable and independently demonstrable, so that the project delivers a working, evaluable result on each even if time constraints limit further extension. Additional capabilities envisioned for a full production platform -- such as a dedicated AML network-analysis agent, a continuously updated customer risk profiler, and an interactive fraud-investigation graph UI -- are discussed as part of the proposed design in Chapter~5, but are treated as stretch goals beyond the three core objectives above.
